\begin{frame}{스케일이 사라지는 것을 직접 계산으로}
  대칭 데이터 예: 최적해 $w=(1/6,1/6),\ b=0$, 서포트 벡터 $(3,3)$.
  \small
  \begin{tabular}{@{}lcc@{}}
    \toprule
     & 기능적 마진 & 기하적 마진 \\
    \midrule
    $w=(1/6,1/6)$, $b=0$ &
      $1\cdot(3/6+3/6)=1$ &
      $\frac{1}{\|(1/6,1/6)\|}=\frac{6}{\sqrt{2}}\approx 4.24$ \\
    \midrule
    $w'=(1/2,1/2)$, $b'=0$ (3배 스케일) &
      $1\cdot(1/2\cdot3+1/2\cdot3)=3$ ($3$배로 커짐) &
      $\frac{3}{\|(1/2,1/2)\|}=\frac{6}{\sqrt{2}}\approx 4.24$ (\textbf{그대로}) \\
    \bottomrule
  \end{tabular}
  \vskip0.8em
  \begin{itemize}
    \item 경계선 $x_1+x_2=0$은 \textbf{아무것도 안 변했는데}
      기능적 마진은 3배로 커졌다.
    \item 실제 거리(기하적 마진)는 스케일에 무관하게 일정.
    \item ``마진 1''이라는 숫자 자체는 관습 — 중요한 것은 \textbf{비율}.
  \end{itemize}
\end{frame}
